%----------------------------------------------------------------------------------------
%	CONFIGURAÇÃO DO DOCUMENTO
%----------------------------------------------------------------------------------------
\documentclass[10pt,a4paper]{article}	
\usepackage[utf8]{inputenc}		
\usepackage[T1]{fontenc}
\usepackage[default]{raleway}
\usepackage[a4paper, margin=1.5cm, top=1.2cm, bottom=1.5cm]{geometry}
\usepackage{xcolor}
\usepackage{hyperref}
\usepackage{fontawesome5}
\usepackage{enumitem}
\usepackage{multicol}

% Definição de Cores
\definecolor{sectcol}{RGB}{0,100,200}
\definecolor{textcol}{RGB}{50,50,50}
\definecolor{lightgray}{RGB}{120,120,120}

\hypersetup{
    colorlinks=true,
    urlcolor=sectcol,
    pdfauthor={Gabriel Matheus Oliveira},
    pdftitle={Currículo - Gabriel Matheus Oliveira}
}

% Ajuste de parágrafo e espaçamento
\setlength{\parindent}{0mm}
\setlist[itemize]{nosep, leftmargin=1em, labelsep=0.5em} % Listas compactas

%----------------------------------------------------------------------------------------
%	COMANDOS CUSTOMIZADOS
%----------------------------------------------------------------------------------------

% Seção estilizada
\newcommand{\cvsection}[1]{
    \vspace{12pt}
    {\large\textcolor{sectcol}{\textbf{\MakeUppercase{#1}}}}
    \vspace{2pt}
    \hrule height 0.5pt \relax
    \vspace{6pt}
}

% Evento de experiência/formação
% Uso: \cvevent{data}{cargo}{empresa}{descrição}
\newcommand{\cvevent}[4]{
    \noindent \textbf{#2} \hfill \textcolor{sectcol}{\small \faCalendar* #1} \\
    \noindent \textit{\textcolor{lightgray}{#3}} \vspace{4pt} \\
    {#4} \vspace{8pt}
}

%----------------------------------------------------------------------------------------
%	INÍCIO DO DOCUMENTO
%----------------------------------------------------------------------------------------
\begin{document}

%--- CABEÇALHO ---
\begin{center}
    {\Huge \textbf{GABRIEL MATHEUS OLIVEIRA}} \\
    \vspace{6pt}
    {\color{sectcol} \large \faCode \ Analista de Sistemas | Técnico em Automação Industrial} \\
    \vspace{10pt}
    
    \small
    \faMapMarker* Uberaba -- MG \quad 
    \faPhone* (34) 99289-1596 \quad 
    \faEnvelope \ \href{mailto:gab_matheus@hotmail.com}{gab\_matheus@hotmail.com} \\
    \vspace{4pt}
    \faLinkedin \ \href{https://linkedin.com/in/oliveiragm}{linkedin.com/in/oliveiragm} \quad 
    \faGithub \ \href{https://github.com/GabMatheus}{github.com/GabMatheus} \quad 
    \faGlobe \ \href{https://gabmatheus.github.io/Portfolio/}{https://gabmatheus.github.io/Portfolio}
\end{center}

%--- RESUMO ---
\cvsection{Resumo Profissional}
Analista de Sistemas com formação superior e técnica. Experiência multidisciplinar que integra o desenvolvimento de software e web, automação de processos corporativos e automação industrial. Construção e desenvolvimento de projetos e scripts em Python para automação de rotinas, suporte especializado a sistemas ERP e gestão de infraestrutura de TI. Atuando também no setor industrial na área de automação e instrumentação, elaborando e modificando projetos elétricos no AutoCad e realizando programação de CLPs.

%--- EXPERIÊNCIA ---
\cvsection{Experiência Profissional}
\cvevent{03/2026 – Atualmente no cargo}{Técnico de Automação}{KR Automação e Robótica}{
\begin{itemize}
    \item Análise de malhas de controle, programação de CLPs Siemens utilizando a plataforma TIA Portal e desenvolvimento de sistemas supervisórios (WinCC/SCADA).
    \item Montagem, cabeamento e interligação de painéis elétricos de comando e automação.
    \item Elaboração de desenhos técnicos, esquemas elétricos e documentação técnica de projetos.
    \item Manutenção e desenvolvimento de aplicações internas de software.
\end{itemize}
}

\cvevent{06/2021 – 10/2025}{Analista de Sistemas}{CONEFI LTDA}{
\begin{itemize}
    \item Desenvolvimento de ferramentas de automação de processos (RPA) e scripts em Python para otimização de fluxos internos e eliminação de tarefas manuais.
    \item Integração de sistemas e manipulação de bancos de dados relacionais (SQL) para garantir a consistência e fluxo de informações corporativas.
    \item Suporte técnico especializado nível N2/N3 aos sistemas ERP (EMSys3, AutoSystem, Domínio, Mastermaq), atuando na resolução de falhas e parametrização.
    \item Gestão, manutenção e expansão da infraestrutura local de redes e servidores do departamento de TI.
    \item Elaboração de relatórios operacionais e apoio técnico na tomada de decisão em projetos de melhoria de sistemas.
    \item Gestão técnica de ativos de tecnologia, englobando o processo de especificação de hardware, compras e vendas do setor.
\end{itemize}
}

\cvevent{2016 – 2019}{Instrumentista Industrial II}{DBR Engenharia/Biano Engenharia}{
\begin{itemize}
    \item Manutenção, calibração, parametrização e diagnóstico de instrumentos.   
    \item Instalação, configuração de inversores de frequência, motores elétricos trifásicos e painéis elétricos de comando.  
    \item Montagem eletromecânica industrial de sistemas automatizados e cabeamento estruturado.  
    \item Leitura, interpretação e execução de projetos elétricos complexos e diagramas de interligação.
\end{itemize}
}

\cvevent{2014 – 2016}{Estágio em Engenharia/TI}{TV Integração (Afiliada Rede Globo)}{
\begin{itemize}
    \item Suporte técnico em sistemas de transmissão e manutenção de infraestrutura de redes.
    \item Operação e suporte a usuários.
\end{itemize}
}

%--- FORMAÇÃO ---
\cvsection{Formação Acadêmica}

\cvevent{06/2026}{Análise e Desenvolvimento de Sistemas}{UFV (inicio) / UNINTER (formação)}{
Matriz curricular focada em Engenharia de Software, Análise de Sistemas, Programação Orientada a Objetos (POO), Arquitetura de Software e Gestão de Projetos de TI e Desenvolvimento web.
}

\cvevent{Concluído}{Técnico em Eletrônica}{CEFET-MG}{
Curso abrange automação industrial, circuitos digitais e analógicos, sistemas embarcados e microcontroladores. Foco prático também em programação de CLPs e desenhos elétricos de comando.
}

%--- HABILIDADES ---
\cvsection{Habilidades Técnicas}

\begin{multicols}{2}
    \textbf{\faIndustry \ Automação \& Sistemas}
    \begin{itemize}
        \item Programação de CLPs
        \item Desenhos elétricos no Autocad
        \item Sistemas Embarcados & Microcontroladores
        \item Instrumentação Industrial
        \item Redes de Comunicação e Protocolos Industriais
    \end{itemize}

    \columnbreak

    \textbf{\faLaptopCode \ Desenvolvimento \& TI}
    \begin{itemize}
        \item Banco de Dados (PostgreSQL, MySQL, SQL Server)
        \item HTML, CSS, JavaScript, Python, C, C++
        \item Versionamento Git e Ambiente Linux
        \item Suporte Técnico Avançado e Administração de ERPs
        \item Engenharia de Software e Princípios SOLID
    \end{itemize}
\end{multicols}

%--- CURSOS E EXTRAS ---
\cvsection{Cursos e Certificações}
\cvevent{2019}{Inglês Geral (Nível B2)}{OSCARS International -- Dublin/Irlanda}{
Intercâmbio comunicativo internacional de 6 meses com foco em fluência e proficiência em ambientes profissionais globais.
}

%--- PUBLICAÇÕES ---
\cvsection{Publicações}
\cvevent{2021}{Artigos Técnicos em Sistemas de Informação}{Produção Acadêmica}{
\begin{itemize}
    \item Artigos focados em modelos de qualidade de software, arquitetura baseada em componentes e funcionamento/aplicações de GPUs.
    \item \faGithub \ \href{https://github.com/GabMatheus/Pesquisas_Latex}{Repositório LaTeX com as publicações}
\end{itemize}
}

%--- REFERÊNCIAS ---
\cvsection{Referências Profissionais}
\begin{itemize}[label=\faUserCheck, leftmargin=1.5em]
    \item \textbf{Bruno Lopes} | Gerente de Operações, Postos Alpa | (34) 98802-3389
    \item \textbf{Thiago Matias} | Coordenador de TI, Conefi | (34) 98848-9918
    \item \textbf{Diego Militão} | Supervisor DBR Engenharia | (34) 99184-9385
\end{itemize}

\vfill
\center{\footnotesize \textcolor{lightgray}{Documento gerado em \LaTeX \ | Disponível para validação via GitHub.}}

\end{document}
